\Needspace{5\baselineskip}\textbf{Problems}

\begin{enumerate}[label=\thesection.\arabic*.,ref=\thesection.\arabic*,start=1,series=book-problems]

\item\label{ex:banach-contraction-picard-lindelof-existence} 应用 Banach 压缩映射原理, 证明 ODE 的 Picard-Lindelöf 存在性定理 (Theorem~\ref{thm:ode-existence}).

\item\label{ex:continuous-bilinear-linear-operator-correspondence} 证明, 从 $X$ 到 $Y$ 的连续双线性映射与从 $X$ 到 $\mathcal{L}(X,Y)$ 的线性映射之间存在自然对应.

\item\label{ex:frechet-derivative-chain-rule} 建立 Fréchet 导数的链式法则.

\item\label{ex:pointwise-composition-frechet-derivative} 验证 Example~\ref{ex:pointwise-composition-frechet-derivative-example} 中的断言.

\item\label{ex:inverse-function-higher-differentiability} 验证 Remark~\ref{rem:inverse-function-higher-differentiability}.

\item\label{ex:sobolev-composition-lemma-details} 补全 Lemma~\ref{lem:sobolev-pointwise-composition-differentiability} 的证明细节.

\item\label{ex:implicit-function-fully-nonlinear-elliptic-boundary-problems} 利用隐函数定理, 得到具有非线性边界条件的完全非线性二阶椭圆 PDE 的存在性结果.

\item\label{ex:sine-composition-ltwo-nowhere-differentiability} 考虑复合映射\[L^2(0,1)\ni u(\cdot)\mapsto\sin(u(\cdot))\in L^2(0,1).\]证明, 这个映射处处不可微.

\end{enumerate}
